% TOP_PROBLEM: {"problemId": "problem.unconditional-separation-of-quantum-and-classical-bounded-error-polynomial-time-bqp-bpp", "canonicalTitle": "Unconditional Separation of Quantum and Classical Bounded-Error Polynomial Time (BQP != BPP)", "releaseRank": 19, "primaryDomain": "quantum_information_computation", "primaryDomainLabel": "Quantum information & computation", "displayStatus": "open", "releaseStatus": "open"}
% !TeX program = lualatex
% Definition and sources only; this file is not an LLM solution attempt.
\documentclass[11pt]{article}
\usepackage[margin=25mm]{geometry}
\usepackage{fontspec}
\setmainfont{Latin Modern Roman}
\usepackage{amsmath,amssymb,mathtools}
\usepackage{unicode-math}
\setmathfont{Latin Modern Math}
\usepackage{xurl,hyperref}
\hypersetup{colorlinks=true,urlcolor=blue}
\setcounter{secnumdepth}{0}
\setlength{\parindent}{0pt}
\setlength{\parskip}{5pt}
\setlength{\emergencystretch}{3em}
\begin{document}
\section*{\texorpdfstring{Unconditional Separation of Quantum and Classical Bounded-Error Polynomial Time (BQP != BPP)}{Unconditional Separation of Quantum and Classical Bounded-Error Polynomial Time (BQP != BPP)}}\label{19}
\subsection{Definitions and mathematical statement}
Both classes decide total binary languages with error at most $1/3$ on every input and polynomial running time. The class $\mathsf{BPP}$ uses classical random bits; $\mathsf{BQP}$ uses uniform quantum computation over a fixed effective universal gate set. The target is
\[
 \exists L\subseteq\{0,1\}^*:\quad L\in\mathsf{BQP}\ \text{and}\ L\notin\mathsf{BPP}.
\]
The assertion is unconditional and unrelativized: a separation only in the presence of an oracle or under a conjectured hardness assumption does not meet it.
\par\smallskip\textit{Formulation and concept references:} \href{https://doi.org/10.1137/S0097539796300921}{[S1]}.

\subsection{Short English statement}
Can one prove, without computational assumptions, that quantum computers efficiently solve some decision problem beyond all efficient randomized classical computers?
\subsection{Sources}
\begin{itemize}
\item[S1]E. Bernstein, U. Vazirani, SIAM J. Comput. 26(5):1411-1473, 1997.
\url{https://doi.org/10.1137/S0097539796300921}.
\textit{Formulation source.}
\item[S2]Improved Separation Between Quantum and Classical Computers for Sampling and Functional Tasks — Marshall, Aaronson and Dunjko, CCC 2025.
\url{https://drops.dagstuhl.de/storage/00lipics/lipics-vol339-ccc2025/LIPIcs.CCC.2025.5/LIPIcs.CCC.2025.5.pdf}.
\textit{Further reference; not a proof endorsement.}
\item[S3]Oracle Separations in the Fourier Hierarchy — Atul Mantri, September 10, 2026.
\url{https://arxiv.org/abs/2609.11830}.
\textit{Further reference; not a proof endorsement.}
\end{itemize}
\end{document}
